%% toc: P15 161 161.10,161.11,161.13
% How Durrett and Gnedenko--Kolmogorov place the normal-domain criterion; domains of attraction of laws and of types; the non-normal stable limit of symmetric power-tail sums.
\chapter{Normal Attraction in Context: Attributions, Types and Stable Limits}
\label{sec:p15-w02c-p15-016}

The criterion of \cref{thm:p15-nda-characterization} settles exactly
when normalized i.i.d.\ sums can have a normal limit. This chapter
records how its two sources place that criterion, and proves what
happens on the excluded side for the power-tail laws of
\cref{thm:p15-nda-examples}: their sums still converge after a
different normalization, but to a limit that is not normal.

\emph{Inputs.} An i.i.d.\ sequence \(X_1,X_2,\ldots\) of finite real
random variables with \(S_n=X_1+\cdots+X_n\); for the stable limit, a
symmetric common law with \(\mathbb P(|X_1|>x)=x^{-\alpha}\) for
\(x\geq1\) and a fixed \(0<\alpha<2\).
\emph{Conclusions.} The domain of attraction depends only on the type
of the attracting law (\cref{thm:p15-w02c-p15-016-types}); the
power-tail sums satisfy \(S_n/n^{1/\alpha}\Rightarrow Y\) with
\(\mathbb Ee^{itY}=\exp(-C|t|^\alpha)\), a nondegenerate law that is
not \(\chi\) and not any normal law
(\cref{thm:p15-w02c-p15-016-stable}); at \(\alpha=1\) the
limit is a scaled Cauchy law
(\cref{thm:p15-w02c-p15-016-cauchy-case}).

The notation \(q_X\), \(H_X\), \(\chi\) and \(N(0,\tau^2)\) is that of
\cref{def:p15-nda-law-and-moments}. Characteristic functions are those
of \cref{sec:p15-characteristic-functions-and-inversion}, whose
algebra, continuity theorem and Gaussian transforms are used without
change.

\section{How the two sources frame the criterion}
\label{sec:p15-w02c-p15-016-framing}

\paragraph{Durrett's framing.}
Durrett states the criterion as his Theorem~3.4.14: for i.i.d.\
\(X_1,X_2,\ldots\), constants \(a_n\) and \(b_n>0\) with
\((S_n-a_n)/b_n\Rightarrow\chi\) exist if and only if
\[
 y^2P(|X_1|>y)/E(|X_1|^2;|X_1|\leq y)\to0
\]
\cite[p.~152]{DurrettPTEAuthor2019}. He attributes the theorem to
L\'evy, and introduces it as the final statement about convergence to
the normal distribution, in his words ``the last word'' on that
question. He does not prove it; he sends the reader to Gnedenko and
Kolmogorov (1954), which he calls ``a reference that contains the
last word on many results about sums of independent random
variables'' \cite[p.~152]{DurrettPTEAuthor2019}. Both directions of the
criterion, in the notation
\(y^2P(|X_1|>y)/E(|X_1|^2;|X_1|\leq y)=y^2q_X(y)/H_X(y)\), are
proved in \cref{thm:p15-nda-characterization}; that statement adds the
hypothesis that the law is nondegenerate, which the printed statement
leaves implicit and which part~4 of \cref{thm:p15-nda-examples} shows
cannot be dropped.

Immediately before the theorem, Durrett remarks that if the tail
\(x^{-2}\) of his Example~3.4.13 is replaced by \(x^{-\alpha}\) with
\(0<\alpha<2\), then \(S_n/n^{1/\alpha}\) converges weakly to a limit
that is not \(\chi\) \cite[p.~152]{DurrettPTEAuthor2019}. The remark
announces the computation for a later section and names it
Section~3.6; in the version cited here the computation is carried out
in Section~3.8, display (3.8.1) to (3.8.3)
\cite[Section~3.8, pp.~182--183]{DurrettPTEAuthor2019}.
\Cref{thm:p15-w02c-p15-016-stable} proves the remark in
full.

\paragraph{The 1954 framing.}
Gnedenko and Kolmogorov define attraction for an arbitrary limit law,
not only the normal one, and only then specialize.

\begin{hbdefinition}{def:p15-w02c-p15-016-attraction}{Attraction of laws and of types}
Let \(\xi_1,\xi_2,\ldots\) be independent with common distribution
function \(F(x)\), and let \(V(x)\) be a distribution function. The
law \(F\) is \emph{attracted} to \(V\) if for suitably chosen
constants \(A_n\) and \(B_n\) the distribution functions of
\[
 \zeta_n=\frac1{B_n}\sum_{k=1}^n\xi_k-A_n
\]
converge to \(V(x)\) as \(n\to\infty\), at every continuity point of
\(V\). The set of all distribution functions attracted to \(V\) is the
\emph{domain of attraction} of \(V\)
\cite[Section~35, pp.~171--172]{GnedenkoKolmogorov1954LimitDistributions}.
Here \(A_n\) is real and \(B_n\) is real and nonzero, so that
\(\zeta_n\) is defined.

Two distribution functions \(F_1\) and \(F_2\) belong to the
\emph{same type} if \(F_2(x)=F_1(ax+b)\) for some constants \(a>0\)
and \(b\); this relation is symmetric and transitive, so the
distribution functions split into disjoint types. All normal laws form
one type; the point masses form the \emph{improper} type, and the
laws of every other type are \emph{proper}
\cite[p.~40]{GnedenkoKolmogorov1954LimitDistributions}.
\end{hbdefinition}

\begin{hblemma}{thm:p15-w02c-p15-016-types}{proved=annex:p15-w02c-p15-016-types}{A domain of attraction depends only on the type}
If \(F\) is attracted to \(V\) with constants \(A_n,B_n\), and
\(W(x)=V(ax+b)\) with \(a>0\) and real \(b\), then \(F\) is attracted
to \(W\) with the constants \(aB_n\) and \((A_n+b)/a\). Consequently
two laws of the same type have the same domain of attraction, and the
domain of attraction of a type is well defined.
\end{hblemma}

The source draws this consequence and then states its programme
\cite[p.~172]{GnedenkoKolmogorov1954LimitDistributions}. It asserts,
by appeal to its \S\,33, that a law has a nonempty domain of
attraction exactly when it is stable. It calls the determination of these
domains one of the fundamental problems in the theory of stable laws
and says that its Section~35 gives the complete solution. It
illustrates the outcome in advance: the normal law attracts a very
wide class of laws, whereas the domain of attraction of each other
stable law contains only laws whose character resembles that of the
attracting law. Its Theorem~1 is then the normal case: \(F\) belongs
to the domain of attraction of a proper normal law if and only if
\[
 \frac{X^2\int_{|x|>X}dF(x)}{\int_{|x|<X}x^2\,dF(x)}\to0
 \qquad(X\to\infty)
\]
\cite[Section~35, Theorem~1, p.~172]{GnedenkoKolmogorov1954LimitDistributions}.
A footnote to the theorem attributes it to A.~Ya.~Khintchine, W.~Feller
and P.~L\'evy, citing items [52], [27] and [75] of the book's
bibliography: Khintchine's 1935 paper on the domain of attraction of
the Gaussian law, Feller's paper on the central limit theorem in
Mathematische Zeitschrift (1935; second part 1937), and L\'evy's 1935
paper on asymptotic properties of sums of independent or chained
random variables
\cite[p.~172, footnote; Bibliography, pp.~258--260]{GnedenkoKolmogorov1954LimitDistributions}.
The source's open truncation \(|x|<X\) gives the same criterion as
the closed truncation of \cref{def:p15-nda-law-and-moments}, by the
endpoint equivalence of \cref{thm:p15-nda-tail-equivalence}.

Within this handbook, the reach of the two framing claims is as
follows. The width of the normal domain is exhibited by
\cref{thm:p15-nda-characterization} and \cref{thm:p15-nda-examples}:
it contains every nondegenerate finite-variance law and also
infinite-variance laws such as the density \(|x|^{-3}\) on
\(|x|>1\). The contrasting claim about the other stable laws is
illustrated, not proved in general, by the next section: the
symmetric power-tail law of index \(\alpha<2\) is attracted to a
non-normal law whose characteristic function
\(\exp(-C|t|^\alpha)\) carries the same index \(\alpha\). The
assertion that only stable laws have nonempty domains, and the
determination of the non-normal domains, are the source's
\S\S\,33 and~35 results and are reported here with that attribution.

\section{Power tails: a non-normal limit at scale \texorpdfstring{\(n^{1/\alpha}\)}{n to the 1/alpha}}
\label{sec:p15-w02c-p15-016-power-tails}

Part~3 of \cref{thm:p15-nda-examples} shows that for the symmetric
law with \(\mathbb P(|X|>x)=x^{-\alpha}\), \(x\geq1\), \(0<\alpha<2\),
no deterministic centering and positive normalization yields a
nondegenerate normal limit. The sums do not thereby fail to converge.

\begin{hbtheorem}{thm:p15-w02c-p15-016-stable}{proved=annex:p15-w02c-p15-016-stable}{Symmetric power tails give a non-normal stable limit}
Let \(0<\alpha<2\), and let \(X_1,X_2,\ldots\) be i.i.d.\ with
\(X_1\) and \(-X_1\) equal in law and
\[
 \mathbb P(|X_1|>x)=x^{-\alpha}\qquad(x\geq1).
 \]
Equivalently \(\mathbb P(X_1>x)=\mathbb P(X_1<-x)=x^{-\alpha}/2\)
for \(x\geq1\), which is Durrett's (3.8.1). Put
\(\varphi(t)=\mathbb Ee^{itX_1}\) and
\[
 C=\alpha\int_0^\infty\frac{1-\cos u}{u^{\alpha+1}}\,du .
 \]
Then:
\begin{enumerate}
\item \(0<C<\infty\), \(\varphi\) is real, and
 \(1-\varphi(t)=\alpha\int_1^\infty(1-\cos tx)x^{-\alpha-1}\,dx\);
\item for every real \(t\),
 \(n\{1-\varphi(t/n^{1/\alpha})\}\to C|t|^\alpha\);
\item \(S_n/n^{1/\alpha}\Rightarrow Y\), where \(Y\) is a random
 variable with \(\mathbb Ee^{itY}=\exp(-C|t|^\alpha)\);
\item the law of \(Y\) is nondegenerate, is not the law of \(\chi\),
 and is not the law \(N(m,v)\) of any normal variable with real mean
 \(m\) and variance \(v\geq0\).
\end{enumerate}
Thus \(S_n/n^{1/\alpha}\) converges weakly to a limit which is not
\(\chi\), as remarked at \cite[p.~152]{DurrettPTEAuthor2019} and
computed at \cite[Section~3.8, pp.~182--183]{DurrettPTEAuthor2019}.
\end{hbtheorem}

\begin{hbcorollary}{thm:p15-w02c-p15-016-cauchy-case}{proved=annex:p15-w02c-p15-016-cauchy-case}{Worked case: index one gives a scaled Cauchy limit}
For \(\alpha=1\), the constant is \(C=\pi/2\), and
\[
 S_n/n\Rightarrow(\pi/2)\,Z,
 \]
where \(Z\) has the standard Cauchy density \(1/(\pi(1+x^2))\).
Here \(\mathbb E|X_1|=\infty\). The sample mean \(S_n/n\) does not
converge in probability to any constant; without further normalization
it has a nondegenerate Cauchy limit.
\end{hbcorollary}

The corollary should be compared with
\cref{thm:p15-characteristic-inverse-examples}: there the summands are
themselves Cauchy and \(S_n/n\) has exactly their law for every \(n\);
here the summands have the two-sided Pareto law and only the limit is
Cauchy.

\begin{hbwarning}{def:p15-w02c-p15-016-power-tail-limits}{What the power-tail theorem does and does not cover}
The theorem uses symmetry: it makes \(\varphi\) real and removes any
centering. It uses the exact tail \(x^{-\alpha}\) above \(1\), which
fixes the law of \(|X_1|\); a law whose tail is only asymptotically
proportional to \(x^{-\alpha}\), or an asymmetric law, is not covered
by this statement. At \(\alpha=2\) the hypothesis becomes Durrett's
Example~3.4.13, the logarithmic case of part~2 of
\cref{thm:p15-nda-examples}: \(\int_0^\infty(1-\cos u)u^{-3}\,du\)
diverges at zero, and the normal limit appears at the scale
\(\sqrt{n\log n}\). Hence \(S_n/\sqrt n\) cannot converge weakly
there: a weak limit of \(S_n/\sqrt n\) would force
\(S_n/\sqrt{n\log n}\Rightarrow0\), contradicting the normal limit. For \(\alpha>2\) the variance is
finite and part~1 of \cref{thm:p15-nda-examples} gives the normal
limit at scale \(\sqrt n\). For \(\alpha<2\) the theorem gives the
non-normal limit, and part~3 of \cref{thm:p15-nda-examples} excludes a
normal one at every normalization.
\end{hbwarning}

\begin{proofboundary}{def:p15-w02c-p15-016-boundary}{Proved and reported claims}
\Cref{thm:p15-w02c-p15-016-types},
\cref{thm:p15-w02c-p15-016-stable} and
\cref{thm:p15-w02c-p15-016-cauchy-case} are proved in the annex from
the characteristic-function results already established. The
attribution of Theorem~3.4.14 to L\'evy and Durrett's description of
it and of the 1954 book are reported from
\cite[p.~152]{DurrettPTEAuthor2019}. The attribution of
Section~35, Theorem~1 to Khintchine, Feller and L\'evy, the claim that
exactly the stable laws have nonempty domains of attraction, the claim
of a complete solution of the domain problem, and the contrast between
the normal domain and the other stable domains are reported from
\cite[p.~172]{GnedenkoKolmogorov1954LimitDistributions}; apart from
the normal case and the power-tail instance proved here, the
non-normal domain theory of that Section~35 and the stable-law
characterization of its \S\,33 are not reproduced in this handbook.
\end{proofboundary}
